\chapter{Resultados de usabilidad}\label{chap:usabilidad}

Este capítulo aplica a los tres frameworks la rúbrica de usabilidad definida en el
Capítulo~\ref{chap:metodologia}, organizada en torno a las cuatro dimensiones de la
experiencia de uso \textit{useful}, \textit{usable}, \textit{desirable} y
\textit{accessible}. A diferencia de la comparativa de rendimiento, que mide el
comportamiento cuantitativo del modelo, y del análisis de flexibilidad, que describe la
capacidad arquitectónica de cada herramienta para adaptarse a escenarios distintos del
previsto, aquí se valora el coste y la experiencia de desarrollo: qué esfuerzo exigió
poner en marcha cada framework, qué fricciones aparecieron durante la implementación y la
depuración, y con qué claridad permitió trabajar. La descripción técnica de cada solución
—arquitectura de archivos, adaptación del experimento y automatización— se detalla en el
Capítulo~\ref{chap:implementacion}; este capítulo no la repite, sino que la interpreta
desde el punto de vista del desarrollador.

La valoración es una evaluación experta cualitativa aplicada por el autor, no un estudio
de experiencia de usuario con participantes. Cada dimensión se puntúa en la escala de uno
a cinco descrita en el Capítulo~\ref{chap:metodologia} y se acompaña de la evidencia que
la justifica: pasos de instalación reales, errores encontrados y resueltos, claridad de
los registros y de los mensajes de error, y barreras de entrada de la documentación. Para
mitigar la subjetividad, los mismos subcriterios y la misma escala se aplican por igual a
los tres frameworks, recogidos sobre un único equipo y un mismo flujo experimental.

\section{Usabilidad de Flower}\label{sec:usab-flower}

\subsection{Utilidad funcional (\textit{useful})}\label{sec:usab-flower-useful}

Flower cubre por completo el caso de uso del estudio. Permitió implementar
\textit{Federated Averaging} sobre CIFAR-10 con una ResNet-18 adaptada, ejecutar varios
clientes mediante simulación, recoger métricas de entrenamiento y evaluación, y reproducir
los experimentos a partir de una configuración declarativa. La separación entre la
aplicación de cliente, la aplicación de servidor y la lógica de tarea encaja de forma
natural con el patrón cliente-servidor del aprendizaje federado, de modo que cada
responsabilidad del experimento tiene un lugar evidente en el proyecto. La única limitación
funcional relevante desde la óptica del uso es que la evaluación federada permanece
desactivada cuando la fracción de evaluación es nula, lo que obliga a apoyar la medición de
la precisión en una evaluación centralizada en el servidor; no es un fallo, sino una
decisión de diseño que conviene conocer de antemano para no interpretar como anómalo el
aviso correspondiente.

\subsection{Facilidad de uso (\textit{usable})}\label{sec:usab-flower-usable}

La puesta en marcha de Flower fue de las más directas del estudio. La instalación se
resolvió con un gestor de paquetes estándar y un conjunto reducido de dependencias, sin
necesidad de recurrir a entornos especiales ni a versiones concretas del intérprete. La
estructura del proyecto es compacta y predecible, y la configuración de cada experimento se
concentra en un único fichero declarativo, lo que reduce el número de lugares donde
intervenir para cambiar el número de rondas, la fracción de clientes o los hiperparámetros.

Durante las primeras ejecuciones aparecieron incidencias de configuración, como la ausencia
de una clave esperada por el experimento, que se manifestaban como errores claros y
localizables. Estas incidencias eran atribuibles a la propia parametrización inicial y no al
framework, y se corrigieron en una sola sesión de trabajo propagando la configuración donde
correspondía. La fricción de uso más significativa apareció en el modo concurrente, basado
en un planificador de procesos que mantiene actores vivos entre rondas y acumula presión de
memoria; en el equipo de pruebas, con memoria limitada, esa acumulación condujo a un agotamiento
de memoria que impidió completar las ejecuciones concurrentes y obligó a tratarlas como
limitación experimental. En el modo secuencial, en cambio, el framework se comportó de forma
estable y predecible.

\subsection{Experiencia de desarrollo (\textit{desirable})}\label{sec:usab-flower-desirable}

Como entorno de trabajo, Flower transmite una sensación de orden y limpieza. Su interfaz de
programación se integra de forma idiomática con el ecosistema de PyTorch, de manera que el
bucle de entrenamiento, la definición del modelo y la carga de datos se escriben de forma
casi idéntica a un entrenamiento no federado, y solo la frontera de comunicación introduce
conceptos específicos del framework. Esa cercanía reduce la carga cognitiva y aumenta la
confianza al modificar el experimento. La contrapartida en la experiencia es que la capa de
ejecución concurrente añade una complejidad que rompe esa sensación de control: cuando el
planificador de procesos interviene, el comportamiento se vuelve menos transparente y la
depuración de los problemas de memoria resulta menos inmediata que en el resto del flujo.

\subsection{Accesibilidad (\textit{accessible})}\label{sec:usab-flower-accessible}

Flower presenta una barrera de entrada baja. Su documentación es abundante y está orientada
a casos prácticos, con ejemplos que sirven de punto de partida directo, y sus mensajes de
registro durante la ejecución resultan legibles y suficientes para seguir la evolución de las
rondas sin instrumentación adicional. Para una persona con conocimientos previos de PyTorch,
el tiempo desde la instalación hasta la primera ejecución útil es corto. Las guías cubren con
claridad la configuración básica, de modo que la mayor parte del esfuerzo inicial se dedica al
experimento en sí y no a entender la herramienta.

\subsection{Valoración de la usabilidad de Flower}\label{sec:usab-flower-valoracion}

En conjunto, Flower es el framework más accesible y agradable de usar de los tres para el
caso de estudio, con un coste de aprendizaje bajo y un flujo secuencial muy estable; su
punto débil de usabilidad se concentra en el modo concurrente. La
Tabla~\ref{tab:usab-flower} resume la valoración por dimensiones.

\begin{table}[htbp]
  \centering
  \caption{Valoración de la usabilidad de Flower según las cuatro dimensiones de la
  rúbrica (escala 1--5).}
  \label{tab:usab-flower}
  \begin{tabular}{llc}
    \toprule
    \textbf{Dimensión} & \textbf{Síntesis de la evidencia} & \textbf{Puntuación} \\
    \midrule
    \textit{Useful}     & Cobertura funcional completa; evaluación federada opcional        & 5 \\
    \textit{Usable}     & Instalación directa; estable en secuencial; concurrencia frágil   & 4 \\
    \textit{Desirable}  & API limpia e idiomática con PyTorch; capa concurrente menos clara & 4 \\
    \textit{Accessible} & Documentación práctica y abundante; logs legibles                 & 5 \\
    \midrule
    \textbf{Media}      &                                                                   & \textbf{4{,}5} \\
    \bottomrule
  \end{tabular}
\end{table}

\section{Usabilidad de NVIDIA FLARE}\label{sec:usab-flare}

\subsection{Utilidad funcional (\textit{useful})}\label{sec:usab-flare-useful}

NVIDIA FLARE cubre el caso de estudio y excede ampliamente sus necesidades. A través de una
plantilla de trabajo que reproduce el patrón de dispersión y agregación equivalente a
\textit{Federated Averaging}, permitió ejecutar el mismo experimento con CIFAR-10, la
ResNet-18 adaptada y varios clientes simulados, registrar métricas por ronda y realizar una
evaluación centralizada final sobre el conjunto de prueba. Su orientación hacia escenarios
de producción aporta funcionalidades muy por encima de lo requerido, lo que en términos de
utilidad es una ventaja, pero también explica que una parte del esfuerzo inicial se dedique
a discernir qué subconjunto de sus capacidades es pertinente para una comparativa académica
controlada.

\subsection{Facilidad de uso (\textit{usable})}\label{sec:usab-flare-usable}

La facilidad de uso es la dimensión más débil de NVIDIA FLARE y la que más diferencia su
experiencia respecto a Flower. La instalación no fue inmediata: el primer intento sobre un
entorno virtual con una versión muy reciente del intérprete falló durante la compilación de
dependencias nativas, y solo se resolvió migrando a un entorno gestionado por Conda con una
versión del intérprete contrastada, sobre la que se instalaron además los extras necesarios.
Esta dependencia de una combinación concreta de versiones constituye ya una primera barrera.

Una vez instalado, el framework exige asimilar un número considerable de conceptos propios
—trabajo, espacio de trabajo, sitio, ejecutor, flujo de trabajo, persistor y agregador,
entre otros— y una estructura de directorios extensa, de modo que la distancia entre instalar
la herramienta y lanzar el experimento previsto es notable. La adaptación requirió resolver
una sucesión de incidencias de naturaleza diversa: dependencias ausentes para la carga de
datos, un agotamiento de memoria del sistema al preparar una prueba con muchos clientes que
llegó a derribar la sesión gráfica del equipo, la selección incorrecta de dispositivo al
invocar rutinas de la GPU cuando el entrenamiento debía realizarse en otro dispositivo, el
cierre prematuro del agente de comunicación cuando el cliente esperaba mensajes de un trabajo
ya finalizado, y errores derivados del parcheo automático de la configuración que llegaban a
alterar el número de rondas o de clientes efectivos respecto al experimento planificado. Cada
una de estas incidencias era resoluble, pero su acumulación sitúa la curva de aprendizaje
claramente por encima de la del resto de herramientas evaluadas.

\subsection{Experiencia de desarrollo (\textit{desirable})}\label{sec:usab-flare-desirable}

La experiencia de desarrollo con NVIDIA FLARE es ambivalente. En su contra juega una
depuración exigente: los fallos de un cliente pueden manifestarse de forma indirecta como
errores de comunicación —pérdida del par, ausencia de latido o aborto de la tarea— que
ocultan la causa real y obligan a inspeccionar varios artefactos antes de localizar el
origen del problema. A su favor, el framework genera un conjunto muy completo de registros y
artefactos por ejecución, con resúmenes, trazas, monitorización del sistema y registros de
evaluación separados, lo que proporciona una trazabilidad experimental superior a la del
resto de herramientas y transmite confianza una vez superada la fase de configuración. El
resultado es una experiencia que recompensa la inversión de aprendizaje con rigor y
reproducibilidad, pero que penaliza con dureza al desarrollador en las primeras fases.

\subsection{Accesibilidad (\textit{accessible})}\label{sec:usab-flare-accessible}

La accesibilidad de NVIDIA FLARE se ve condicionada por la profundidad del framework. La
documentación es extensa y existe abundante material de referencia, pero la barrera de
entrada es alta: empezar a producir resultados requiere un conocimiento previo no trivial de
su modelo de ejecución, y algunos comandos o vías esperadas no coincidían con las
efectivamente disponibles en la versión instalada, lo que obliga a contrastar la
documentación con la herramienta real. Los mensajes de error, especialmente los relacionados
con la comunicación entre procesos, no siempre son autoexplicativos para quien se inicia.
Para un perfil con experiencia en sistemas distribuidos la documentación es valiosa, pero
para una primera aproximación la facilidad de acceso es limitada.

\subsection{Valoración de la usabilidad de NVIDIA FLARE}\label{sec:usab-flare-valoracion}

NVIDIA FLARE es el framework más potente y trazable, pero también el más costoso de usar:
combina una utilidad funcional sobresaliente con una facilidad de uso y una accesibilidad
penalizadas por su complejidad y su curva de aprendizaje. La Tabla~\ref{tab:usab-flare}
resume la valoración.

\begin{table}[htbp]
  \centering
  \caption{Valoración de la usabilidad de NVIDIA FLARE según las cuatro dimensiones de la
  rúbrica (escala 1--5).}
  \label{tab:usab-flare}
  \begin{tabular}{llc}
    \toprule
    \textbf{Dimensión} & \textbf{Síntesis de la evidencia} & \textbf{Puntuación} \\
    \midrule
    \textit{Useful}     & Cobertura completa y capacidades muy por encima de lo requerido       & 5 \\
    \textit{Usable}     & Instalación sensible a versiones; muchos conceptos; depuración costosa & 2 \\
    \textit{Desirable}  & Trazabilidad excelente, pero depuración indirecta y exigente          & 3 \\
    \textit{Accessible} & Documentación amplia con barrera de entrada alta; errores opacos      & 3 \\
    \midrule
    \textbf{Media}      &                                                                       & \textbf{3{,}25} \\
    \bottomrule
  \end{tabular}
\end{table}

\section{Usabilidad de FedML}\label{sec:usab-fedml}

\subsection{Utilidad funcional (\textit{useful})}\label{sec:usab-fedml-useful}

FedML cubre el caso de estudio y aporta una característica útil desde el punto de vista del
uso: la disponibilidad de distintos backends de ejecución, uno para el modo secuencial y otro
basado en paso de mensajes para el modo concurrente, bajo una misma configuración. Permitió
ejecutar \textit{Federated Averaging} sobre CIFAR-10 con la ResNet-18 adaptada y los mismos
hiperparámetros, así como obtener un \textit{baseline} repetido con varias semillas. La
principal limitación funcional en términos de uso es que el modo concurrente no produjo una
métrica de precisión final aprovechable en sus registros, lo que restringe su utilidad a la
medición de tiempo y recursos y obliga a tenerlo en cuenta al planificar qué resultados puede
sustentar cada modo.

\subsection{Facilidad de uso (\textit{usable})}\label{sec:usab-fedml-usable}

La facilidad de uso de FedML se sitúa en un punto intermedio. La instalación y la puesta en
marcha no presentaron fricciones reseñables, y la configuración se gestiona mediante ficheros
declarativos que centralizan los parámetros del experimento. La fricción se concentró en dos
aspectos. El primero es la fiabilidad de la instrumentación: en varias ejecuciones la
inicialización de la biblioteca de monitorización de la GPU falló por un desajuste entre la
versión del controlador y la de la biblioteca, de modo que, aunque las métricas de
aprendizaje seguían siendo válidas, las métricas de GPU de esas corridas no eran utilizables
y debían marcarse como no disponibles en lugar de mezclarse con ceros. El segundo es la
extracción de resultados: en el modo concurrente, los registros no siempre contenían la
información necesaria de forma directa, y algunas ejecuciones terminaron por agotamiento de
tiempo o sin recoger por completo los resultados, lo que exigió un trabajo de filtrado más
cuidadoso que en Flower.

\subsection{Experiencia de desarrollo (\textit{desirable})}\label{sec:usab-fedml-desirable}

Como entorno de desarrollo, FedML resulta funcional y razonablemente cómodo. El mecanismo de
extensión por herencia, que separa el entrenador del cliente y el agregador del servidor,
ofrece puntos de intervención claros y una sensación de orden al adaptar el experimento. La
experiencia se ve algo mermada por la menor pulcritud de los registros en determinados modos
y por la necesidad de validar con cuidado qué corridas son aprovechables, lo que introduce una
capa de verificación adicional sobre la salida del framework. El modo concurrente, basado en
paso de mensajes, se mostró en algunas pruebas más estable que el de Flower, lo que mejora la
percepción de robustez en ese escenario.

\subsection{Accesibilidad (\textit{accessible})}\label{sec:usab-fedml-accessible}

La accesibilidad de FedML es intermedia. Dispone de documentación y de ejemplos que permiten
arrancar, pero parte de la configuración y, sobre todo, la interpretación y extracción de las
métricas requieren un esfuerzo de lectura mayor que en Flower, y algunos comportamientos solo
se comprenden tras inspeccionar los registros con detenimiento. No impone barreras de entrada
tan altas como NVIDIA FLARE, pero tampoco alcanza la inmediatez de Flower para pasar de la
instalación a un resultado limpio y bien documentado.

\subsection{Valoración de la usabilidad de FedML}\label{sec:usab-fedml-valoracion}

FedML ofrece una usabilidad correcta y equilibrada, sin la fricción de instalación de NVIDIA
FLARE pero sin la inmediatez de Flower; sus puntos débiles son la fiabilidad de la
instrumentación de GPU y la claridad de los registros en el modo concurrente. La
Tabla~\ref{tab:usab-fedml} resume la valoración.

\begin{table}[htbp]
  \centering
  \caption{Valoración de la usabilidad de FedML según las cuatro dimensiones de la rúbrica
  (escala 1--5).}
  \label{tab:usab-fedml}
  \begin{tabular}{llc}
    \toprule
    \textbf{Dimensión} & \textbf{Síntesis de la evidencia} & \textbf{Puntuación} \\
    \midrule
    \textit{Useful}     & Cobertura completa; doble backend; concurrente sin precisión final & 4 \\
    \textit{Usable}     & Instalación sin fricción; instrumentación de GPU y logs irregulares & 3 \\
    \textit{Desirable}  & Extensión por herencia clara; verificación extra de resultados     & 3 \\
    \textit{Accessible} & Documentación suficiente; extracción de métricas exigente          & 3 \\
    \midrule
    \textbf{Media}      &                                                                    & \textbf{3{,}25} \\
    \bottomrule
  \end{tabular}
\end{table}

\section{Comparativa de usabilidad}\label{sec:usab-comparativa}

La Tabla~\ref{tab:usab-comparativa} reúne las doce puntuaciones de las cuatro dimensiones
para los tres frameworks, lo que permite una lectura transversal de la usabilidad.

\begin{table}[htbp]
  \centering
  \caption{Comparativa de la valoración de usabilidad de los tres frameworks según las
  cuatro dimensiones de la rúbrica (escala 1--5).}
  \label{tab:usab-comparativa}
  \begin{tabular}{lcccc}
    \toprule
    \textbf{Dimensión} & \textbf{Flower} & \textbf{NVIDIA FLARE} & \textbf{FedML} \\
    \midrule
    \textit{Useful}     & 5 & 5 & 4 \\
    \textit{Usable}     & 4 & 2 & 3 \\
    \textit{Desirable}  & 4 & 3 & 3 \\
    \textit{Accessible} & 5 & 3 & 3 \\
    \midrule
    \textbf{Media}      & \textbf{4{,}5} & \textbf{3{,}25} & \textbf{3{,}25} \\
    \bottomrule
  \end{tabular}
\end{table}

La lectura por dimensiones revela un patrón coherente. En utilidad funcional los tres
frameworks obtienen valoraciones altas, ya que todos permiten cubrir el caso de estudio; las
diferencias proceden de limitaciones puntuales, como la ausencia de precisión final en el
modo concurrente de FedML. Es en la facilidad de uso donde la distancia es máxima: Flower
ofrece una puesta en marcha directa y un flujo secuencial estable, FedML se sitúa en un punto
intermedio condicionado por la fiabilidad de su instrumentación, y NVIDIA FLARE queda por
detrás por su instalación sensible a versiones, su gran número de conceptos y su depuración
costosa. En accesibilidad se reproduce la misma ordenación, con Flower claramente por delante
gracias a una barrera de entrada baja y unos registros legibles. La dimensión de experiencia
de desarrollo es la más matizada, porque la excelente trazabilidad de NVIDIA FLARE compensa
parcialmente su dificultad y eleva su valoración por encima de lo que sugeriría su facilidad
de uso aislada.

De esta lectura se desprende una recomendación de uso dependiente del perfil del proyecto.
Para prototipado rápido, docencia o trabajos en los que el tiempo desde la instalación hasta
el primer resultado es crítico, Flower es la opción más usable. Para proyectos que prioricen
la trazabilidad, la reproducibilidad rigurosa y la cercanía a un despliegue real, y que puedan
asumir una curva de aprendizaje alta, NVIDIA FLARE compensa su coste de uso con la calidad de
sus artefactos. FedML ocupa una posición intermedia adecuada cuando interesa disponer de
varios backends de ejecución bajo una misma configuración, siempre que se asuma un trabajo
adicional de verificación y extracción de métricas.

\section{Síntesis}\label{sec:usab-sintesis}

La evaluación de usabilidad responde al objetivo de comparar la experiencia práctica de uso
de cada framework y confirma que la elección de una herramienta de aprendizaje federado no
puede basarse únicamente en su rendimiento. Los tres frameworks resultaron funcionalmente
aptos para el caso de estudio, pero ofrecieron experiencias de desarrollo marcadamente
distintas: Flower destacó por su accesibilidad y su bajo coste de aprendizaje, NVIDIA FLARE
por su trazabilidad a costa de una notable complejidad, y FedML por un equilibrio razonable
con limitaciones concretas en la instrumentación y la extracción de resultados. Conviene
recordar que estas valoraciones reflejan el juicio del autor como desarrollador sobre un único
entorno de pruebas, por lo que deben interpretarse como una guía cualitativa y no como una
medida absoluta; esta condición se recoge entre las amenazas a la validez del estudio. El
capítulo siguiente analiza la flexibilidad de cada framework, un eje complementario que
describe hasta qué punto cada herramienta permite adaptarse a escenarios distintos del
previsto por defecto.